\documentclass{article}

\begin{document}
\begin{quotation}
When you restart the server (kill the running server and run make run-server
again), what happens to subsequent client lookups (get operations)? How does the
client know that the server restarted?
\end{quotation}
If the previous tree was not empty, then the subsequent client lookups will fail
with \texttt{Exception("Root hash mismatch").}

If the previous tree was empty, then the lookups will succeed, but return None.

The client knows because it will store the previous root hash in merkle.root
before exiting and read it upon starting.

\begin{quotation}
Why do subsequent client updates (put operations) fail to execute after a server
restart? What would go wrong if they didn’t fail?
\end{quotation}
Because the server was not in the state (represented by the root hash) that the
client expected the server to be.

What would go wrong if they didn't fail? Well, then it would be as if there was
no check at all.


\begin{quotation}
How many siblings are included in the proof to verify the node for get foo after
you’ve inserted both foo=bar and hello=world into the store, as in the above
example? Explain why that is.
\end{quotation}
Nine siblings. It turns out that they have the first eight bits of their hashes
equal.


\end{document}
